\documentclass[11pt,twoside]{amsart}
\usepackage{amssymb, amsmath, enumerate, palatino, hyperref}
\usepackage[normalem]{ulem}
\usepackage{fullpage}
\usepackage[T1]{fontenc}
\renewcommand{\labelitemi}{\guillemotright}
\usepackage{mathrsfs}

\usepackage{tikz}
\usetikzlibrary{positioning}

\tikzstyle{ball} = [circle,shading=ball, ball color=black,
    minimum size=1mm,inner sep=1.3pt]


\theoremstyle{plain}
\newtheorem{prop}{Proposition}%[section]
\newtheorem{lemma}[prop]{Lemma}
\newtheorem{thm}[prop]{Theorem}
\newtheorem{obs}[prop]{Observation}
\newtheorem{app}[prop]{Application}
\newtheorem*{MainThm}{Main Theorem}
\newtheorem{cor}[prop]{Corollary}
\newtheorem{conj}[prop]{Conjecture}
\theoremstyle{remark}
\newtheorem{rmk}[prop]{Remark}
\newtheorem{prob}[prop]{Problem}
\newtheorem{q}[prop]{Question}
\newtheorem{bonus}[prop]{Bonus Problem}
\newtheorem{exc}{Exercise}
\theoremstyle{definition}
\newtheorem{ex}[prop]{Example}
\theoremstyle{definition}
\newtheorem{defn}[prop]{Definition}

\newcommand{\RR}{\mathbb{R}}
\newcommand{\ZZ}{\mathbb{Z}}
\newcommand{\CC}{\mathbb{C}}
\newcommand{\NN}{\mathbb{N}}
\newcommand{\QQ}{\mathbb{Q}}
\newcommand{\PP}{\mathbb{P}}

\newcommand{\cC}{\mathscr{C}}
\newcommand{\Ob}{\operatorname{Ob}}
\newcommand{\Mor}{\operatorname{Mor}}

\newcommand{\Set}{\operatorname{Set}}
\newcommand{\FinSet}{\operatorname{FinSet}}
\newcommand{\Vect}{\operatorname{Vect}}
\newcommand{\FinVect}{\operatorname{FinVect}}
\newcommand{\Mat}{\operatorname{Mat}}
\newcommand{\Gp}{\operatorname{Gp}}
\newcommand{\FinGp}{\operatorname{FinGp}}
\newcommand{\AbGp}{\operatorname{AbGp}}
\newcommand{\Ring}{\operatorname{Ring}}
\newcommand{\CommRing}{\operatorname{CommRing}}
\newcommand{\Field}{\operatorname{Field}}
\newcommand{\Top}{\operatorname{Top}}

\newcommand{\Aut}{\operatorname{Aut}}

\newcommand{\defeq}{:=}

\title{Math 113: Discrete Structures\\ Homework 25}

\begin{document}
\maketitle

\noindent
\textbf{Due:} Monday, April 6 at 10pm.
\bigskip

\begin{prob}\label{independence} 
Suppose a fair coin is flipped~$n>1$ times. We record the result at a string of
length~$n$ in the letters~$H$ and~$T$.  For instance, if~$n=3$, then~$HTT$ says
that the first flip was heads and the next two were tails.  Our sample
space~$S$ is, thus, the set of strings of length~$n$ consisting of the
letters~$H$ and~$T$, upon which we place the uniform distribution.
Let~$A$ be the event that the first flip is heads, and let~$B$ be the event that
the last flip is heads.  

\begin{enumerate}[(a)]
    \item Prove that these events are independent by computing
 the relevant probabilities and using the definition of independence.
   \item For the case of $n=2$, give an example of a probability distribution on $S$ for which these same two events $A$ and $B$ are \emph{not} independent. Note that to give a probability distribution you will need to assign probabilities to each of the outcomes, i.e., to each of the elements of $S$.
\end{enumerate}
\end{prob}


\begin{prob}\label{birthday problem} This is the famous birthday problem.
  Suppose there are~$n$ people in a room.  For simplicity, we will say that
  there are~$365$ possible birthdays (i.e., we will ignore leap years) and that
  each day is equally likely to be a birthday. To create the sample space,~$S$,
  number the people from~$1$ to~$n$, then the possible outcomes are the
  sequences of length~$n$ where the~$i$-th element of the sequence is a possible
  birthday for person~$i$. Let~$B$ be the event that at least two people in the
  room share a birthday.
  \begin{enumerate}[(a)]
    \item What is~$|S|$?
    \item Consider the complementary event~$B^c$, i.e., the event that no two
      people have the same birthday.  Give an expression for~$B^c$, and use it
      to find the probability~$P(B^c)$.  (The result will depend on~$n$.)
    \item Use the previous result to give an expression for~$P(B)$.
    \item A plot of~$P(B)$ as a function
      of~$n$ looks like this:
      \begin{center}
	\includegraphics[height=1.5in]{birthday.png}
      \end{center}
      Use a calculator of some sort to give decimal approximations, accurate to
      three decimal places, for~$P(B)$ when~$n=22$ and~$n=23$. You should see
      that if the room contains at least~$23$ people, then the probability two
      people share the same birthday is more than one half.
  \end{enumerate}
\end{prob}
\end{document}
